\section{Introduction}

For a positive integer \(r\), let
\(\operatorname{rad}(r)=\prod_{p\mid r}p\). We count ordered pairs
\[
 F(X)=\#\{(m,n)\in\mathbb Z_{>0}^2:m<n\le X,\
       \operatorname{rad}(m(m+1))=\operatorname{rad}(n(n+1))\}.
\]

\begin{theorem}\label{thm:main}
For every \(\varepsilon>0\) and \(X\ge2\),
\[
                         F(X)\ll_\varepsilon X^{109/165+\varepsilon}.
\]
\end{theorem}

The exponent is \(1/165\) below \(2/3\). Section~\ref{sec:notation}
introduces the four-corner gcd coordinates. It also records the
source-specific identities used here. The new part has three pieces: a
first-match partition into thirteen disjoint populations, uniform
estimates for its new rows, and a finite optimization of the long-shift
remainder.

The minimum \(S_0\) below is the limiting minimum of the row savings.
After the analytic loss \(\beta_0\) and a conservative long-row charge
\(g_0/4\),
the proof retains the smaller saving displayed below.

The earlier result \cite[Theorem~1]{Paper1} gives exponent \(2/3\) for
this counting function. Its support-fibre lemma and low-radical estimate
are used below; Appendix~\ref{app:sources} states them and records the
substitutions. The quotient-tail and Farey-spacing estimates in that paper
are context only. They are not used. Two related papers by the author
\cite{Zhang712,ZhangE20} count primitive \(abc\) triples, not the pairs
counted here, and do not enter this proof.

At the chosen rational point, the minimum row saving is
\[
 S_0=\frac{30869}{5000000},\qquad
 g_0=S_0-\frac1{165}=\frac{18677}{165000000}.
\]
The source-uniform estimates spend a total analytic loss
\(\beta_0=g_0/16\). The long row has three finite corrections: geometric,
interval, and high-degree. Beyond the common threshold each is at most
\(g_0/4\). We charge their maximum, not their sum. The remaining margin is
\[
 S_0-\beta_0-\frac{g_0}{4}
 =\frac1{165}+\frac{11g_0}{16}>\frac1{165}.
\]

Sections~\ref{sec:notation}--\ref{sec:partition} give the canonical source
and the exact row partition. Sections~\ref{sec:outer}--\ref{sec:intermediate}
treat the outer, joint-corner, short-shift and intermediate rows.
Section~\ref{sec:long} handles the long row. Sections~\ref{sec:ledger}--\ref{sec:conclusion}
allocate losses and sum the dyadic shells. The appendices contain the
self-contained square-fibre proof, the application of the cited estimates, and the
branch arithmetic. The use of AI tools is described after the appendices.

We use Heath--Brown's published Lemma~2 under its stated bihomogeneous
hypotheses. The fixed cubic branch uses the unnumbered theorem of
Bombieri--Schmidt with its 1989 correction \cite{BS1989}. Appendix~\ref{app:sources}
states both inputs and records the applications.
Reuss's Theorem~1 \cite{Reuss} is the determinant-method input behind the low-radical
estimate in \cite{Paper1}. That estimate is used here. Tao's
Lemma~2.10 gives a comparison with the square-fibre estimate of
Appendix~\ref{app:fibre}. The proof uses that appendix theorem.
